\section{The empirical distribution}

A dataset can be turned into a probability distribution of its own. This does not
mean that the dataset becomes the true population law. It means that we define a
probability measure that gives equal weight to each observed data point.

Suppose the data are
\[
x_1,x_2,\ldots,x_n.
\]
The empirical distribution assigns mass \(1/n\) to each observation. If the same
value occurs several times, the masses at that value add together.

\begin{definitionbox}
For a dataset \(x_1,\ldots,x_n\), the \textbf{empirical distribution} \(P_n\) is
defined by
\[
P_n(A)=\frac1n\sum_{i=1}^n \mathbf{1}_{\{x_i\in A\}}
\]
for every set \(A\) of values under consideration.
\end{definitionbox}

The indicator \(\mathbf{1}_{\{x_i\in A\}}\) equals \(1\) if the observation \(x_i\)
lies in \(A\), and equals \(0\) otherwise. Therefore the sum counts how many
observations fall in \(A\). Dividing by \(n\) gives the observed proportion.

\subsection*{Example: repeated values}

Consider the dataset
\[
2,\quad 3,\quad 2,\quad 5,\quad 3,\quad 3,\quad 4,\quad 2.
\]
The empirical distribution assigns mass \(1/8\) to each observation. Since the
value \(2\) appears three times, its total empirical mass is
\[
P_n(\{2\})=\frac38.
\]
Similarly,
\[
P_n(\{3\})=\frac38,
\qquad
P_n(\{4\})=\frac18,
\qquad
P_n(\{5\})=\frac18.
\]
The empirical distribution is therefore a discrete probability law concentrated
on the observed values.

\subsection*{A continuous sample still gives a discrete empirical distribution}

Suppose the data are measurements from a continuous process:
\[
1.82,\quad 2.17,\quad 1.96,\quad 2.04,\quad 2.31.
\]
Even if the theoretical model is continuous, the empirical distribution of this
finite dataset is discrete. It puts mass \(1/5\) at each of the five observed
numbers.

This is an important point. A continuous model may generate data, but once we
observe a finite sample, the empirical distribution is a finite discrete object.
The continuity belongs to the model, not to the finite list of observed values.

\begin{remarkbox}
Every finite empirical distribution is discrete. This remains true even when the
data are measurements from a continuous phenomenon.
\end{remarkbox}

\subsection*{Relation to the theoretical distribution}

If data are generated from a random variable \(X\) with theoretical probability
law \(P\), then \(P\) describes the mechanism before observation. The empirical
distribution \(P_n\) describes the observed sample after observation.

For a set \(A\), compare
\[
P(X\in A)
\]
with
\[
P_n(A)=\frac{\#\{i:x_i\in A\}}{n}.
\]
The first is a model probability. The second is an observed proportion. If the
model is correct and the sample is large, the observed proportion may be close to
the model probability. But closeness is not automatic for every finite sample.

\subsection*{Empirical distribution as a model of the data}

The empirical distribution is a perfect probability model for the observed data
in the following limited sense: if one chooses one of the recorded observations
uniformly at random, then the probability of choosing a value in \(A\) is exactly
\(P_n(A)\).

For example, if the dataset has \(100\) observations and \(23\) of them lie in the
interval \([0,1]\), then choosing one observation uniformly at random from the
recorded dataset gives probability
\[
\frac{23}{100}
\]
of selecting an observation in \([0,1]\).

This interpretation is useful, but it is not the same as saying that the true
population probability of \([0,1]\) is \(23/100\). The empirical distribution is
about the sample.

\begin{summarybox}
The empirical distribution assigns mass \(1/n\) to each observation. It is the
probability law of a randomly selected observed data point. It summarizes the
sample exactly, while a theoretical distribution describes an idealized random
mechanism that may have generated the sample.
\end{summarybox}
